% 《大学物理简明教程》第5章解答题5.3.1—5.3.9，印刷/PDF第125—126页。
% 参数从两页扫描图片逐项核对，原创解析，待独立学术复核。
\subsection{解答题：分子动理论的定量应用}

\sourceproblem{5.3.1}{125}{125}
\problemtext 容器中的氧气压强$p=0.1\,\mathrm{MPa}$、温度$27^\circ\mathrm C$。
求分子数密度$n$、单个氧分子质量$m$和气体质量密度$\rho$。
\workings
采用$T=300.15$ K、$p=10^5$ Pa、$M_{\mathrm{O_2}}=0.028\,\mathrm{kg/mol}$：
\[
\boxed{n=\frac{p}{k_{\mathrm B}T}\approx
2.41\times10^{25}\,\mathrm{m^{-3}}},
\quad
\boxed{m=\frac{M_{\mathrm{O_2}}}{N_A}
\approx4.65\times10^{-26}\,\mathrm{kg}},
\]
\[
\boxed{\rho=nm=\frac{pM_{\mathrm{O_2}}}{RT}
\approx1.12\,\mathrm{kg/m^3}}.
\]
\crosscheck 将$n$乘单分子质量所得密度与理想气体状态方程求得的一致。

\sourceproblem{5.3.2}{125}{125}
\problemtext 一水银气压计原水银柱高$0.76\,\mathrm m$，玻璃管顶端至水银液面的空隙为$0.12\,\mathrm m$，
管横截面积$2.0\times10^{-4}\,\mathrm{m^2}$。注入少量He后水银柱高降至$0.60\,\mathrm m$。
温度$27^\circ\mathrm C$、水银密度$1.36\times10^4\,\mathrm{kg/m^3}$，求氦气质量（摩尔质量$0.004\,\mathrm{kg/mol}$）。
\workings
忽略原真空残余气体压强，注入后He的压力等于水银柱减少部分所对应的压差：
\[
p_{\mathrm{He}}=\rho_{\mathrm{Hg}}g(0.76-0.60)
=2.13248\times10^4\,\mathrm{Pa}.
\]
此时上方气体所占长度增加到
$0.12+(0.76-0.60)=0.28$ m，故
$V=(0.28)(2.0\times10^{-4})=5.6\times10^{-5}\,\mathrm{m^3}$。
\[
\boxed{m_{\mathrm{He}}
=\frac{p_{\mathrm{He}}VM_{\mathrm{He}}}{RT}
\approx1.91\times10^{-6}\,\mathrm{kg}}.
\]
\crosscheck 压差乘气体体积具有能量量纲；必须使用\emph{下降后的气柱高度}计算气体体积，不能沿用初始0.12 m。

\sourceproblem{5.3.3}{125}{125}
\problemtext 理想气体分子方均根速率为$400\,\mathrm{m/s}$，
压强为1 atm。求气体质量密度。
\workings
压强和方均根速率满足
$p=\frac13\rho\langle v^2\rangle$，故
\[
\boxed{\rho=\frac{3p}{v_{\mathrm{rms}}^2}
=\frac{3(101325)}{400^2}
\approx1.90\,\mathrm{kg/m^3}}.
\]
\crosscheck 压强若取$1.013\times10^5$ Pa，结果的有效位数不变。

\sourceproblem{5.3.4}{125--126}{125--126}
\problemtext 一瓶氧气和一瓶氢气处于相同温度；
已知氢气分子的平均平动动能
$\overline\varepsilon_{\mathrm H_2}=6.21\times10^{-21}\,\mathrm J$。
求氧气分子平均平动动能及方均根速率，并求氢气温度。
\workings
分子平均平动动能只由绝对温度决定，
故氧气分子也为$\boxed{6.21\times10^{-21}\,\mathrm J}$。
\[
T=\frac{2\overline\varepsilon}{3k_{\mathrm B}}
\approx\boxed{299.9\,\mathrm K}.
\]
氧气方均根速率：
\[
\boxed{v_{\mathrm{rms,O_2}}
=\sqrt{\frac{3RT}{0.032\,\mathrm{kg/mol}}}
\approx483\,\mathrm{m/s}}.
\]
\crosscheck 同温氢气的方均根速率约为氧气4倍，不能因两者平动平均动能相同就认为平均速率相同。

\sourceproblem{5.3.5}{126}{126}
\problemtext 外太空每立方厘米中约有一个氢气分子，
温度$T=3.5\,\mathrm K$。求氢气分子的方均根速率和气体压强。
\workings
数密度$n=1\,\mathrm{cm^{-3}}=10^6\,\mathrm{m^{-3}}$。
以氢气分子质量$m_{\mathrm H_2}\approx2(1.66054\times10^{-27})\,\mathrm{kg}$，
\[
\boxed{v_{\mathrm{rms}}
=\sqrt{\frac{3k_{\mathrm B}T}{m_{\mathrm H_2}}}
\approx209\,\mathrm{m/s}},
\]
\[
\boxed{p=nk_{\mathrm B}T\approx4.83\times10^{-17}\,\mathrm{Pa}}.
\]
\crosscheck 这是低密度经典理想气体的形式估算；极稀薄宇宙空间的真实气体不必处于局部热平衡。

\sourceproblem{5.3.6}{126}{126}
\problemtext 体积$V=1\,\mathrm{m^3}$，其中含
$N_1=1.0\times10^{25}$个氧气分子和$N_2=4.0\times10^{25}$个氢气分子，
混合气体温度$400\,\mathrm K$。求分子的总平动动能及混合气体压强。
\workings
总分子数$N=N_1+N_2=5.0\times10^{25}$。由自由度均分定理与道尔顿定律：
\[
\boxed{E_{\mathrm{trans}}=\frac32Nk_{\mathrm B}T
\approx4.14\times10^5\,\mathrm J},
\quad
\boxed{p=\frac{Nk_{\mathrm B}T}{V}
\approx2.76\times10^5\,\mathrm{Pa}}.
\]
\crosscheck 即便两种气体分子质量不同，平均平动动能仍由共同温度确定；两种气体分压相加为总压。

\sourceproblem{5.3.7}{126}{126}
\problemtext 解释分子动理论中的六种能量表达式：
$\frac12k_{\mathrm B}T$、$\frac32k_{\mathrm B}T$、
$\frac i2k_{\mathrm B}T$、
$\frac{M}{M_{\mathrm{mol}}}\frac i2RT$、
$\frac i2RT$和$\frac32RT$的物理含义。
\workings
能量均分定理下，每个二次型自由度的平均能量为
$\boxed{\frac12k_{\mathrm B}T}$；
一个分子的三维平动动能均值为
$\boxed{\frac32k_{\mathrm B}T}$；
一个具有$i$个有效自由度的分子的平均总能量为
$\boxed{\frac i2k_{\mathrm B}T}$。
质量为$M$、摩尔质量$M_{\mathrm{mol}}$的理想气体有
$n=M/M_{\mathrm{mol}}$ mol，其内能为
$\boxed{U=n\frac i2RT}$。
$\frac i2RT$是1 mol同类气体的内能；
$\frac32RT$是1 mol任意经典理想气体的\emph{平动}总动能。
\crosscheck $k_{\mathrm B}T$是每分子尺度、$RT$是每摩尔尺度，$R=N_Ak_{\mathrm B}$，勿混用。

\sourceproblem{5.3.8}{126}{126}
\problemtext 容器内混合了11 kg二氧化碳和2 kg氢气，
题设将\emph{两种气体均近似为刚性双原子分子理想气体}；
已知混合气体内能$U=8.1\times10^6$ J。
求温度以及两种气体每分子的平均动能。
\workings
依\emph{题设有效自由度$i=5$}（尽管真实CO$_2$是线性三原子分子），
$n_1=11/0.044=250$ mol、$n_2=2/0.002=1000$ mol，
总摩尔数1250 mol。
\[
U=\frac52(n_1+n_2)RT
\quad\Longrightarrow\quad
\boxed{T=\frac{2U}{5(1250)R}\approx312\,\mathrm K}.
\]
两种分子的平均总动能相同，均为
\[
\boxed{\overline\varepsilon
=\frac52k_{\mathrm B}T\approx1.08\times10^{-20}\,\mathrm J}.
\]
\crosscheck 特别依赖题目要求的“五自由度”近似；真实CO$_2$的自由度激发程度依温度而异，此处不能擅自用真实分子结构替代教材题设。

\sourceproblem{5.3.9}{126}{126}
\problemtext 温度$27^\circ\mathrm C$时1 mol氢气分子的
平动动能、转动动能及内能分别是多少（刚性双原子模型）？
\workings
$T=300.15$ K，故
\[
\boxed{E_{\mathrm{trans}}=\frac32RT\approx3.74\times10^3\,\mathrm J},
\quad
\boxed{E_{\mathrm{rot}}=RT\approx2.50\times10^3\,\mathrm J},
\]
\[
\boxed{U=\frac52RT\approx6.24\times10^3\,\mathrm J}.
\]
\crosscheck 上述三者满足$E_{\mathrm{trans}}+E_{\mathrm{rot}}=U$。
